\documentclass[11pt]{article}
\usepackage[margin=25mm]{geometry}
\usepackage{fontspec}
\setmainfont{Latin Modern Roman}
\usepackage{amsmath,amssymb,mathtools}
\usepackage{unicode-math}
\setmathfont{Latin Modern Math}
\usepackage{xurl,hyperref}
\hypersetup{colorlinks=true,urlcolor=blue}
\setcounter{secnumdepth}{0}
\setlength{\parindent}{0pt}
\setlength{\parskip}{5pt}
\setlength{\emergencystretch}{3em}
\newcommand{\sref}[2]{\hyperlink{source:#1:#2}{[S#2]}}
\newcommand{\cataloguescope}[1]{\paragraph{Recorded target.}#1}
\begin{document}
\section{Multiple ergodic averages — Problem 24}\label{statement}

\subsection{Definitions and mathematical statement}

\paragraph{Shared notation.}$\mathbb N=\{1,2,\ldots\}$, and $\mathbb Z,\mathbb Q,\mathbb R,\mathbb C$ denote the integers, rationals, reals and complex numbers. Any other conventions are specified locally.


A system is a probability space $(X,\mathcal X,\mu)$ with invertible measure-preserving transformations; multiple transformations are assumed to commute. Write $T^j f=f\circ T^j$, and $T^j A$ for the image of a measurable set under $T^j$. A system with one transformation is ergodic if every invariant measurable set has measure zero or one. The space $L^\infty(\mu)$ consists of essentially bounded complex measurable functions. Choose distinct $a,b\in(0,\infty)\setminus\mathbb Z$. Here $\lfloor t\rfloor$ denotes the largest integer at most $t$.
\paragraph{Statement recorded in the supplied source.}
Does every ergodic system and every pair $f,g\in L^\infty(\mu)$ satisfy \[\lim_{N\to\infty}\frac1N\sum_{n=1}^N f(T^{\lfloor n^a\rfloor}x)g(T^{\lfloor n^b\rfloor}x)=\left(\int_X f\,d\mu\right)\left(\int_X g\,d\mu\right)\] for $\mu$-almost every $x$? \sref{statement}{2}
\subsection{Short English statement}
Do two distinct positive noninteger powers give almost-everywhere convergence to the product of means in every ergodic system?
\subsection{Sources}
\begin{description}
\item[\hypertarget{source:statement:1}{S1}] ULAM.AI and the UnsolvedMath Contributors, \emph{UnsolvedMath: A Curated Collection of Open Mathematics Problems}, catalogue entry. CC BY 4.0. \url{https://huggingface.co/datasets/ulamai/UnsolvedMath}.
\item[\hypertarget{source:statement:2}{S2}] N. Frantzikinakis, Some open problems on multiple ergodic averages (2016), Section 6.1.3, full Problem 24, p.32. \url{https://arxiv.org/pdf/1103.3808v3}.
\end{description}
\label{end:statement}
\end{document}
